\documentclass[10pt,a4paper]{amsart}
\usepackage[margin=19mm]{geometry}
\usepackage{amsmath,amssymb,mathtools}
\usepackage[scr=rsfs,cal=euler]{mathalfa}
\usepackage{xfrac}
\usepackage[unicode,hidelinks,bookmarks=false]{hyperref}
\newcommand{\ie}{i.e.\ }
\newcommand{\cf}{cf.\ }
\newcommand{\resp}{respectively\ }
\newcommand{\Subsection}{\subsection}
\providecommand{\nobreakdash}{\nobreak\hbox{-}}
\setlength{\emergencystretch}{2em}
\multlinegap0pt
\usepackage{amssymb}
\usepackage[shortlabels]{enumitem}
\usepackage{tensor}
\usepackage{mathtools}
\usepackage{tikz}
\newtheorem{proposition}{Proposition}
\newcommand\q{\mathbf q}
\renewcommand{\to}{%
   \ifbool{@display}{\longrightarrow}{\rightarrow}%
   }
\title{Cocenter of Hecke algebras of Kac-Moody groups}
\author{Xuhua He, Felix Schremmer}
\date{Source: arXiv:2609.12385v1 (CC BY 4.0)}
\begin{document}
\maketitle

\noindent\emph{Source and licence.} Cocenter of Hecke algebras of Kac-Moody groups. Version arXiv:2609.12385v1 (CC BY 4.0), distributed by arXiv under the Creative Commons Attribution 4.0 International licence, http://creativecommons.org/licenses/by/4.0/, as recorded in the arXiv licence metadata for this exact version and shown on the abstract page https://arxiv.org/abs/2609.12385v1. Full licence notice: \texttt{licenses/arxiv-2609.12385-CC-BY-4.0.txt}.\par\smallskip

\noindent\textit{Editorial scope.} Counted unit: the article's proposition prop:Mbmodule (source0.tex lines 1175--1182). The article's complete original proof of it is delivered with it (source0.tex lines 1183--1232, one \texttt{proof} environment of 50 lines). No mathematics is written, repaired or re-derived: the statement and the proof are the article's own lines, in the article's own notation, and no retained line is substituted or re-flowed.  No further source-local context is needed: every cross-reference of every retained line resolves inside the two delivered spans.  Nothing was omitted from any retained span, so the package carries no omission record.  No retained line contains a citation, so the wrapper carries no bibliography and no bibliography entry.  No retained line contains a figure environment, an image-inclusion command or drawing code, so no image is delivered and the package declares no asset dependency.  Editorial formatting only, in addition: an AMS \texttt{amsart} preamble that restates the article's own declarations and macros used by the retained lines, namely the environments \texttt{proposition} on separate counters, together with the standard packages those lines need.  The retained environments print in this document as Proposition 1 in order of appearance, whereas the article prints the same items with the numbering of its own full document.  The complete native source of the article as distributed in the arXiv e-print for this exact version is bundled unchanged as \texttt{sources/210112/source0.tex}.
\begin{proposition}
\label{prop:Mbmodule}
Let $b$ be a straight element. Set $C_H(T_b)= \{h\in H\mid h T_b = T_b h\}$. Let $M^b$ be the free $\mathbb Z[\q^{\pm 1}]$-module with basis $\{M^b_z\mid z\in W\}$.
For integers $n\geq 1$, we let $f_n:M^b\to H$ be the unique $\mathbb Z[\q^{\pm 1}]$-linear map sending $M^b_z$ to $T_{zb^n}$ for all $z\in W$. Then there exists a unique $H$-$C_H(T_b)$-bimodule structure of $M^b_z$ such that for all $h\in H, m\in M$ and $c\in C_H(T_b)$, there exists some constant $N\geq 1$ with
        \begin{align*}
            f_n(hmc) = h f_n(m) c\text{ for all }n\geq N.
        \end{align*}
\end{proposition}

\begin{proof}
We first show the following auxiliary claim:
For every $m\in M$, there exists some constant $N\geq 1$ with
        \begin{align}\label{eq:Mbmoduleproof}
            f_{n_2}(m) = f_{n_1}(m) T_{b^{n_2-n_1}}\text{ for all }n_2\geq n_1\geq N.
        \end{align}
By linearity, it suffices to consider the case $m = M^b_z$ for some $z\in W$. Consider the sequence
    \begin{align*}
        a_n = \ell(zb^n) - \ell(b^n) = \ell(zb^n) -n\ell(b),~n\geq 1.
    \end{align*}
    It is bounded from below by $-\ell(z)$. It is non-increasing since
    \begin{align*}
        a_{n+1} - a_n = \ell(zb^{n+1}) - \ell(zb^n) - \ell(b) \leq 0.
    \end{align*}
    Hence the sequence eventually stabilizes. Let $N\gg 0$ with $a_N = a_{N+1} = a_{N+2}=\cdots$. Then for $n_2\geq n_1\geq N$, we get
    \begin{align*}
        \ell(zb^{n_2}) = a_{n_2} + n_2 \ell(b) = a_{n_1} + n_2\ell(b) = \ell(zb^{n_1}) + (n_2-n_1)\ell(b) = \ell(zb^{n_1}) + \ell(b^{n_2-n_1}).
    \end{align*}
    Hence, the product $(zb^{n_1})(b^{n_2-n_1})$ is length additive. We conclude
    \begin{align*}
        f_{n_2}(m) = T_{zb^{n_2}} = T_{zb^{n_1}} T_{b^{n_2-n_1}} = f_{n_1}(m) T_{b^{n_2-n_1}}.
    \end{align*}
    Our auxiliary claim \eqref{eq:Mbmoduleproof} is proved.
    
    Note that each map $f_n$ is a $\mathbb Z[\q^{\pm 1}]$-linear endomorphism. We first claim that for each $h\in H, m\in M$ and $c\in C_H(T_b)$, the sequence
    \begin{align*}
        m_n := f_n^{-1}(hf_n(m) c)\in M^b,~n\geq 1
    \end{align*}
    stabilizes: By \eqref{eq:Mbmoduleproof}, we find $N_1\geq 1$ with $f_n(m) = f_{N_1}(m) T_{b^{n-N_1}}$ for all $n\geq N_1$. Set now $h_2 := h f_{N_1}(m) c$ and let $L\geq 1$ be a constant such that
    \begin{align*}
        h_2 \in \sum_{\substack{z\in W\\\ell(z)\leq L}} \mathbb Z[\q^{\pm 1}]T_z.
    \end{align*}
    Then for $n\geq 0$, we get
    \begin{align*}
        h_2 T_{b^n} \in \sum_{\substack{z\in W\\\ell(z)\leq L}} \mathbb Z[\q^{\pm 1}]T_{zb^n} = f_n(M_2),
    \end{align*}
    where $M_2 := \sum_{\substack{z\in W\\\ell(z)\leq L}} \mathbb Z[\q^{\pm 1}]M^b_z$. By \eqref{eq:Mbmoduleproof}, we can find $N_2\geq 1$ such that for all $m\in M_2$ and $n\geq N_2$, we get
    \begin{align*}
        f_n(m) = f_{N_2}(m) T_{b^{n-N_2}}.
    \end{align*}
    Let $m_2\in M_2$ with $h_2 T_{b^{N_2}} = f_{N_2}(m_2)$. Then for all $n\geq N_1+N_2$, we get
    \begin{align*}
        &f_n^{-1}(h f_n(m) c) = f_n^{-1}(h f_{N_1}(m) (T_b)^{n-N_1} c) = f_n^{-1}(h f_{N_1}(m) c (T_b)^{n-N_1})
        \\=&f_n^{-1}(h_2 T_{b^{N_2}} T_{b^{n-N_1-N_2}}) = f_n^{-1}(f_{N_2}(m_2) T_{b^{n-N_1-N_2}})
        =f_n^{-1}(f_{n-N_1}(m_2)).
    \end{align*}
    Using the definition of the $f_\bullet$, it is easy to see that this sequence stabilizes.

    We now define our bimodule structure on $M^b$: For $m\in M, h\in H$ and $c\in C_H(T_b)$, we let $hm$ be the limit value of the sequence $f_n^{-1}(hf_n(m))$ for $n\to \infty$, and $mc$ the limit value of $f_n^{-1}(f_n(m)c)$. The defining properties of a $H$-$C_H(T_b)$-bimodule are now easy to verify, using that $H$ is naturally a $H$-$C_H(T_b)$-bimodule.
\end{proof}

\end{document}
